\documentclass[a4paper,9pt]{article}
\usepackage[utf8]{inputenc}
\usepackage[T1]{fontenc}
\usepackage{booktabs}
\usepackage{enumitem}
\usepackage{amsmath}
\usepackage{amssymb}
\usepackage{listings}
\usepackage{xcolor}
\usepackage{multicol}
\usepackage[
top=0cm,
bottom=1.5cm,
left=0.5cm,
right=0.5cm
]{geometry}


\lstset{
    basicstyle=\ttfamily\small,
    backgroundcolor=\color{gray!10},
    frame=single
}

\title{C Cheat Sheet – Zusammenfassung}
\author{}
\date{}

\begin{document}
    %\maketitle
    %\tableofcontents
    
    \section{C Variable Name Rules (Quick Summary)}
    \begin{multicols}{2}
        \subsection*{Allowed}
        \begin{itemize}
            \item Letters: \texttt{a--z}, \texttt{A--Z}
            \item Digits: \texttt{0--9} (not as first character)
            \item Underscore: \texttt{\_}
            \item Dollar sign: \texttt{\$} (allowed, but not portable)
        \end{itemize}
        
        \subsection*{Not Allowed}
        \begin{itemize}
            \item Starting with a digit (\texttt{1value})
            \item Spaces (\texttt{total sum})
            \item Special characters (\texttt{value@x}, \texttt{my-var})
            \item C keywords (\texttt{int}, \texttt{return}, \texttt{while})
            \item Empty names (\texttt{int ;})
        \end{itemize}
        
        \subsection*{Compiles but Bad Practice}
        \begin{itemize}
            \item Names starting with \texttt{\_}
            \item Very short or confusing names (\texttt{l}, \texttt{O})
        \end{itemize}
        
        \subsection*{Examples}
        \begin{lstlisting}
        int count;
        int total_sum;
        int value2;
        int _temp;     // legal but discouraged
        int $VOLUME;   // legal but not portable
        \end{lstlisting}
    \end{multicols}
    
    \section{printf Format Specifiers}
    \begin{multicols}{2}
    \subsection*{Characters \& Strings}
    \begin{tabular}{lll}
        \toprule
        Specifier & Expects type & Meaning \\
        \midrule
        \texttt{\%c} & int & Single character \\
        \texttt{\%s} & char* & Null-terminated string \\
        \bottomrule
    \end{tabular}
    
    \subsection*{Integers}
    \begin{tabular}{lll}
        \toprule
        Specifier & Expects type & Meaning \\
        \midrule
        \texttt{\%d}, \texttt{\%i} & int & Signed integer \\
        \texttt{\%u} & unsigned int & Unsigned decimal \\
        \texttt{\%o} & unsigned int & Octal \\
        \texttt{\%x}, \texttt{\%X} & unsigned int & Hexadecimal \\
        \bottomrule
    \end{tabular}
    
    \subsection*{Long Integers}
    \begin{tabular}{ll}
        \toprule
        Specifier & Expects type \\
        \midrule
        \texttt{\%ld} & long int \\
        \texttt{\%lu} & unsigned long int \\
        \texttt{\%lld} & long long int \\
        \texttt{\%llu} & unsigned long long int \\
        \bottomrule
    \end{tabular}
    
    \subsection*{Floating Point}
    \begin{tabular}{lll}
        \toprule
        Specifier & Expects type & Meaning \\
        \midrule
        \texttt{\%f}, \texttt{\%lf} & double & Fixed-point \\
        \texttt{\%e} & double & Scientific notation \\
        \texttt{\%g} & double & Shortest of \%f / \%e \\
        \bottomrule
    \end{tabular}
    
    \subsection*{Pointers \& Misc}
    \begin{itemize}
        \item \texttt{\%p} – pointer (\texttt{void*})
        \item \texttt{\%\%} – print \%
        \item \texttt{\%n} – store printed char count (dangerous!)
    \end{itemize}
    
    \subsection*{Critical Rules}
    \begin{itemize}
        \item \texttt{\%s} requires a valid null-terminated string
        \item Wrong specifier $\Rightarrow$ undefined behavior
        \item \texttt{char}, \texttt{short} $\rightarrow$ promoted to \texttt{int}
        \item \texttt{float} $\rightarrow$ promoted to \texttt{double}
    \end{itemize}
    \end{multicols}
    
    \section{C Bitwise Operators}
    
    \begin{tabular}{lll}
        \toprule
        Operator & Name & Description \\
        \midrule
        \texttt{\&} & AND & Bit = 1 if both bits are 1 \\
        \texttt{|} & OR & Bit = 1 if either bit is 1 \\
        \texttt{\^{}} & XOR & Bit = 1 if bits differ \\
        \texttt{\~{}} & NOT & Invert all bits \\
        \texttt{<<} & Left shift & Shift bits left \\
        \texttt{>>} & Right shift & Shift bits right \\
        \bottomrule
    \end{tabular}
    
    \subsection*{Example}
    \begin{lstlisting}
        int a = 6; // 110
        int b = 3; // 011
        a & b; // 2
        a | b; // 7
        a ^ b; // 5
    \end{lstlisting}
    
    \section{Pointer Arithmetic}
    
    \subsection*{Fundamental Rule}
    \[
    p + n = \text{address} + n \cdot \text{sizeof}(T)
    \]
    
    \begin{multicols}{2}
        \subsection*{Valid Operations}
        \begin{itemize}
            \item \texttt{p++}, \texttt{p--}
            \item \texttt{p + n}, \texttt{p - n}
            \item \texttt{p2 - p1} (same array only)
        \end{itemize}
        
        \subsection*{Invalid Operations}
        \begin{itemize}
            \item \texttt{p + q}
            \item \texttt{p * 2}, \texttt{p / 2}
        \end{itemize}
        
        \subsection*{Arrays \& Pointers}
        \begin{lstlisting}
int a[5];
int *p = a;   // same as &a[0]
p[i] == *(p + i);
        \end{lstlisting}
        
        \subsection*{const Pointers}
        \begin{tabular}{|c|c|}
            \hline
            const int *p & Wert konstant \\
            \hline
            int *const p & Pointer konstant \\
            \hline
            const inr *const p & beides konstant \\
            \hline
        \end{tabular}
        
        \subsection{Pointer Precedence}
        \begin{lstlisting}
*p++ // *(p++) p is increased
(*p)++ // value is increased
*p + 1   // Value + 1
*(p + 1) // next element
        \end{lstlisting}
        
        \subsection{Effect of Type}
        \begin{lstlisting}
int *p;     p + 1   // +sizeof(int) (4 bytes)  
double *d;  d + 1   // +sizeof(double) (8 bytes)  
char *c;    c + 1   // +1 byte  
        \end{lstlisting}
        
        \subsection{Pointer Difference}
        \begin{lstlisting}
int a[10];  
int *p1 = &a[2];  
int *p2 = &a[7];  
p2 - p1   // 5   
        \end{lstlisting}
        
        \subsection*{Golden Rule}
        \begin{center}
            \textbf{Pointer arithmetic moves in ELEMENTS, not BYTES.}\\
            \textbf{Pointers may only be dereferenced if the point on to valid memory}
        \end{center}
    \end{multicols}
    
    \section{Calculations with different Datatypes}
    
    \subsection{base rule}
    smaler type becomes bigger type
    \begin{lstlisting}
        char -> short -> int -> long -> float -> double -> long double 
    \end{lstlisting}
    
    \subsection{examples}
    \begin{multicols}{2}
        \subsubsection{intieger calculation}
        \begin{lstlisting}
    int a = 5, b = 2;
    int c = a / b;   // 2
    float d = c; // 2.0
        \end{lstlisting}
        \subsubsection{int/float calculation}
        \begin{lstlisting}
    int a = 5;
    float b = 2;
    float c = a / b;     // 2.5
        \end{lstlisting}
        \subsubsection{cast/float calculation}
        \begin{lstlisting}
    float x = 5.0 / 2;   // 2.5
    float x = 5 / 2.0;   // 2.5
    float x = (float)5/2; // 2.5
        \end{lstlisting}
        \subsubsection{char calculation}
        \begin{lstlisting}
    char c = 'A';
    int x = c + 1;       //66
    printf("%c\n", 'A' + 2); //C
        \end{lstlisting}
        \subsubsection{modulo calculation}
        \begin{lstlisting}
    5 % 2;        // 1
    // 5.0 % 2;  // Error
        \end{lstlisting}
        \subsubsection{printf Output}
        \begin{lstlisting}
    printf("%f\n", 5 / 2);   // 2.000000
    printf("%d\n", (int)5.9); // 5
        \end{lstlisting}
    \end{multicols}
    
    \subsection{flaotingpoints}
    
    \begin{multicols}{2}
        \subsection{Grundregel}
        Gleitkommazahlen ohne Suffix sind `double`
        \begin{lstlisting}
            3.0      // double
            3.0f     // float
            3.0F     // float
        \end{lstlisting}
        \subsection{Gleitkomma-Literals}
        Gleitkommazahlen ohne Suffix sind `double`
        \begin{lstlisting}
            9.1
            9.1f
            7.4e2
            7.4e2f
            1e-3
        \end{lstlisting}
        nicht erlaubt!
        \begin{lstlisting}
    3.5f2     // f darf nicht in der Mitte stehen
    7.4ef     // e braucht eine Zahl
    e2        // incomplete
        \end{lstlisting}
        \subsection{Implizite Typenumwandlung}
        erlaubt
        \begin{lstlisting}
            float  -> double  
            int    -> double   
            int    -> float    
        \end{lstlisting}
        problematisch
        \begin{lstlisting}
        double -> float 
        \end{lstlisting}
        \subsection{Suffixe}
        \begin{tabular}{|c|c|}
            \hline
            keins & double \\
            \hline
            f/F & float \\
            \hline
            e & Exponentialsschreibweise \\
            \hline
            e2 & $10^2$ (double)\\
            \hline
            e-3f & $10^{-3}$ (float) \\
            \hline
        \end{tabular}\\
        Example
        \begin{lstlisting}
            7.4e2    // double (740)
            7.4e2f   // float  (740)    
        \end{lstlisting}
    \end{multicols}
    \subsection{Vergleiche von Fliesskommzahlen}
    \begin{multicols}{2}
        \begin{lstlisting}
            double a = 0.1 + 0.2;
            double b = 0.3;
            double epsilon = 1e-9
            
            if (fabs(a - b) < epsilon)
        \end{lstlisting}
        \begin{tabular}{|c|c|}
            \hline
            typ & epsilon \\
            \hline
            float & 1e-6 \\
            \hline
            double & 1e-9 oder 1e-12 \\
            \hline
        \end{tabular}\\
    \end{multicols}
    \section{fibonacci}
    \begin{lstlisting}
        int fibonacci(int num)
        {
            if (num <= 1) {
                return num;
            }
            return fibonacci(num-2)+fibonacci(num-1);
        }
    \end{lstlisting}
    \pagebreak
    \subsection{arrays}
    \begin{multicols}{2}
        \textbf{initliasisierung}
        \begin{lstlisting}
int x[5] = {1,2,3};    // Rest automatisch 0
int y[]  = {4,5,6};   // Size = 3
        \end{lstlisting}
        \textbf{Zugriff}
        \begin{lstlisting}
a[i]      // lesen
a[i] = 7; // schreiben
        \end{lstlisting}
        \textbf{Size}
        \begin{lstlisting}
sizeof(a)            // gesamte Bytes
sizeof(a)/sizeof(a[0])  // Anzahl Elemente
        \end{lstlisting}
        \textbf{Functiona argument}
        \begin{lstlisting}
void f(int a[], int n);
        \end{lstlisting}
    \end{multicols}
    \subsection{struct}
             \textbf{Struct}
    \begin{lstlisting}
        struct Point {
            int x;
            int y;
        };
    \end{lstlisting}
    \textbf{Deklaration und Verwendung}
    \begin{lstlisting}
        struct Point p;
        p.x = 3;
        p.y = 4;
    \end{lstlisting}
    \textbf{Initialisierung}
    \begin{lstlisting}
        struct Point p1 = {3, 4}; // Reihenfolge wichtig
        struct Point p2 = {.y = 4, .x = 3}; // Designated initializer
    \end{lstlisting}
    \textbf{typedef}
    \begin{lstlisting}
        typedef struct {
            double x;
            double y;
        } Point;
        
        Point p;    // no "struct" needed
    \end{lstlisting}
    \textbf{Zugriff: Pointer}
    \begin{lstlisting}
        p.x        // Access direkt
        ptr->x     // Access via Pointer
        (*ptr).x   // exakt gleich wie ptr->x
    \end{lstlisting}
    \begin{lstlisting}
        ptr->x  =  (*ptr).x
    \end{lstlisting}
        \textbf{Strcut über Pointer verändern}
    \begin{lstlisting}
        void move(Point *p) {
            p->x += 1;
            p->y += 1;
        }
    \end{lstlisting}
    \begin{lstlisting}
        move(&p);
    \end{lstlisting}
   \textbf{Übergabe an Funktionen}
   \begin{lstlisting}
       by value (Kopie)
       
       void f(Point p);
   \end{lstlisting}
   \begin{lstlisting}
       by reference (empfohlen)
       
       void f(Point *p);
   \end{lstlisting}
   
   \section{speicher}
    \begin{lstlisting}
        int *ptr = (int *)malloc(sizeof(int) * 5);
        // Populate the array
        for (int i = 0; i < 5; i++)
        ptr[i] = i + 1;
    \end{lstlisting}
    
\end{document}
